\section{Theory: Formal Results}\label{sec:theory_formal}

This section states the model's main implications in proposition form. The core novelty of the framework is that foreign-currency demand is decomposed into a hedge component and a run component. This decomposition allows the model to generate monotone foreign-currency demand, endogenous public information, crisis amplification, and welfare reversal within a single tractable structure.

\subsection{Setup}

Total foreign-currency demand is

Q_FX = Q_hedge + Q_run.

Hedge demand is given by

Q_hedge = max(0, chi_h * theta - eta_h * p_FX),

where theta is exchange-rate overvaluation and p_FX is the effective acquisition cost of foreign currency. Run demand is

Q_run = gamma_r * (1 + psi * omega_s) * Lambda(pi),

where pi is crisis probability, omega_s is the weight on the public signal, and Lambda is the logistic map.

Crisis probability is

pi = Lambda(Q_FX - (alpha - delta * theta)).

Stablecoin usage improves public information through

sigma_s^2 = barsigma_s^2 / (1 + rho * Q_S),

so that a larger stablecoin market increases the weight placed on the public signal.

\subsection{Propositions}

\paragraph{Proposition 1. Monotone foreign-currency demand.}
Suppose the cost schedule is not so steep that the increase in p_FX dominates the increase in theta. Then hedge demand is weakly increasing in theta. If crisis probability is also weakly increasing in theta, then total foreign-currency demand is weakly increasing in theta.

\emph{Intuition.} The hedge motive ensures that higher exchange-rate misalignment directly raises desired FX holdings even absent panic. The run motive then reinforces this increase once crisis beliefs begin to rise.

\paragraph{Proposition 2. Stablecoin demand rises with overvaluation.}
If stablecoins have lower access costs than cash and the stablecoin premium does not rise too quickly with quantity, then stablecoin demand is weakly increasing in theta.

\emph{Intuition.} As total foreign-currency demand rises with overvaluation, the lower-cost instrument captures a growing share of that demand.

\paragraph{Proposition 3. Endogenous public information.}
Public-signal precision is increasing in stablecoin demand. As a result, the weight placed on the public signal is increasing in stablecoin demand.

\emph{Intuition.} Stablecoin adoption does not merely respond to information; it produces information by making the stablecoin price more informative.

\paragraph{Proposition 4. Information-participation feedback.}
There is a positive feedback between stablecoin participation and public information. Greater stablecoin use improves public-signal precision, which increases the weight placed on public information and thereby raises the coordination component of FX demand.

\paragraph{Proposition 5. Stablecoins increase fragility.}
Holding fixed the underlying overvaluation state, the presence of stablecoins weakly increases crisis probability whenever stablecoins reduce FX access costs and improve public information.

\emph{Intuition.} Stablecoins raise total FX demand both by lowering the effective cost of hedging and by increasing the run component through stronger common information.

\paragraph{Proposition 6. Welfare reversal.}
There exists a threshold level of overvaluation such that stablecoins improve welfare below the threshold but reduce welfare above it.

\emph{Intuition.} At low levels of misalignment, the access benefit dominates. At high levels of misalignment, stablecoins amplify coordinated runs and increase expected disruption losses.

\subsection{Discussion}

These propositions clarify the model's conceptual contribution. Stablecoins are not simply a cheaper vehicle for acquiring foreign currency. They also generate a public signal that reshapes coordination in the parallel market. This dual role explains why stablecoins can simultaneously improve access and increase fragility, and why their welfare effects depend on the underlying degree of exchange-rate misalignment.
